\exercisesection{1}

\begin{enumerate}
\def\labelenumi{\arabic{enumi}.}
\tightlist
\item
  设 A 是 Z 型分次交换环，\((A_{n})\) 是其分次结构；置 \(A^{\geq} = \bigoplus_{n \geq 0} A_{n}, A^{\leq} = \bigoplus_{n \leqslant 0} A_{n}\)；它们是 A 的分次子环。

\begin{enumerate}
\def\labelenumii{(\alph{enumii})}
\item
  对于每个齐次元 \(f\)，它属于 \(\mathbf{A}\)，次数为 \(d\)，分式环 \(\mathbf{A}_f\)（对应于由 \(f^n\)（其中 \(n \geqslant 0\)）组成的乘性子集）具有典范的分次环结构（《代数》，第 II 章，§11）；记 \(\mathbf{A}_{(f)}\) 为 \(\mathbf{A}_f\) 中由次数 0 的元素组成的子环。证明：若 \(d > 0\)，则 \((\mathbf{A}^{\geq})_f = \mathbf{A}_f, \mathbf{A}_{(f)}\) 同构于 \(\mathbf{A}^{(d)} / (f - 1)\mathbf{A}^{(d)}\)，且 \(((\mathbf{A}_f)^{\geq})_{f/1}\) 是与 \(\mathbf{A}_f\) 同构的分次环。
\item
  令多项式环 \(\mathbf{B} = \mathbf{A}[\mathbf{X}] = \mathbf{A} \otimes_{\mathbf{Z}} \mathbf{Z}[\mathbf{X}]\) 具有由 \(\mathbf{A}\) 和 \(\mathbf{Z}[\mathbf{X}]\) 上的分次结构的张量积给出的分次；这个分次与 \(\mathbf{B}\) 的环结构相容。证明：若 \(d > 0\)，则 \(\mathrm{B}_{(f)}\) 同构于 \((\mathrm{A}_f)^{\leqslant}\)，而 \((\mathrm{A}^{(d)})_f\) 是与 \(\mathrm{A}_{(f)} \otimes_{\mathbf{Z}} \mathbf{Z}[\mathbf{X}, \mathbf{X}^{-1}]\) 同构的分次环。
\item
  令 \(g\) 是 \(A\) 的另一个次数为 \(e\) 的齐次元。证明：若 \(d > 0\) 且 \(e > 0\)，则 \(A_{(fg)}\) 同构于 \((A_{(f)})_{g^d / f^e}\)。\\
  \item[* (d)] 假设 \(d > 0\) 且 \(A_n = \{0\}\) 对 \(n < 0\) 成立。证明：若 \(A\) 是 Noether 环，则 \(A_{(f)}\) 也是 Noether 环（使用 §2，第 10 节，定理 2 的推论 4）。*
\end{enumerate}

\item
  设 A 是 Z 型分次交换环，且 \(A_{n}=0\) 对 n\textless0 成立。对于所有 \(n\geqslant0\)，令 \(A_{[n]}\) 表示 A 的分次理想 \(\bigoplus_{m\geqslant n}A_{m}\)；A-代数 \(A^{\sharp}=\bigoplus_{n\geqslant0}A_{[n]}\) 是一个按 \(A_{[n]}=A_{n}^{\sharp}\) 分次的环；对于所有 \(f\in A_{d}(d>0)\)，令 \(f^{\sharp}\) 表示 \(A^{\sharp}\) 中的元素，除次数 d 的分量等于 f 外，其余分量均为零。

\begin{enumerate}
\def\labelenumii{(\alph{enumii})}
\item
  证明：若 \(f \in \mathbf{A}_d\) 且 \(d > 0\)，则环 \(\mathrm{A}_{(f^{\sharp})}^{\sharp}\) 同构于 \((\mathbf{A}_f)^{\geqslant}\)（记号见习题 1）。
\item
  \(\mathbf{A}^{\natural}\) 是有限生成 A-代数的充要条件是 \(\mathbf{A}\) 是有限生成 \(\mathbf{A}_0\)-代数。
\item
  要使 \(\mathbf{A}_{n+1}^{\natural} = \mathbf{A}_1^{\natural}\mathbf{A}_n^{\natural}\)（分别使 \(\mathbf{A}_n^{\natural} = (\mathbf{A}_1^{\natural})^n\)）在 \(n \geqslant n_0\) 时成立，其充要条件是 \(\mathbf{A}_{n+1} = \mathbf{A}_1\mathbf{A}_n\)（分别 \(\mathbf{A}_n = (\mathbf{A}_1)^n\)）在 \(n \geqslant n_0\) 时成立。
\end{enumerate}

\item
  设 K 是一个域，B 是按总次数分次的多项式环 \(\mathrm{K}[X,Y]\)，C 是 B 中由 X 和 \(Y^{2}\) 生成的分次子环，a 是 C 中由 \(Y^{4}\) 生成的分次理想。证明：在分次环 A = C/a 中，\(A_{n+1} = A_{1}A_{n}\) 对 \(n \geqslant 2\) 成立，但 \(A_{n} \neq (A_{1})^{n}\) 对 \(n \geqslant 6\) 成立。
\end{enumerate}
% semantic-fragments: legacy-input-seam
\relax{}
